% Generated draft: requires mathematical and source review.
\documentclass[11pt]{article}
\usepackage[margin=25mm]{geometry}
\usepackage{fontspec}
\setmainfont{Latin Modern Roman}
\usepackage{amsmath,amssymb,mathtools,unicode-math}
\setmathfont{Latin Modern Math}
\usepackage{enumitem,needspace,xurl,hyperref}
\setlength{\parindent}{0pt}
\setlength{\parskip}{5pt}
\setlength{\emergencystretch}{3em}
\newcommand{\N}{\mathbb{N}}
\newcommand{\Z}{\mathbb{Z}}
\newcommand{\Q}{\mathbb{Q}}
\newcommand{\R}{\mathbb{R}}
\newcommand{\C}{\mathbb{C}}
\newcommand{\F}{\mathbb{F}}
\newcommand{\A}{\mathbb{A}}
\newcommand{\PP}{\mathbb{P}}
\newcommand{\E}{\mathbb{E}}
\newcommand{\eps}{\varepsilon}
\newcommand{\dd}{\,\mathrm{d}}
\begin{document}
\section*{UnsolvedMath ID: 30002436}
\textbf{Source problem code:} catalogue entry
\subsection{Definitions and mathematical statement}
Let \(\mathbb{P}\) denote the set of positive prime integers, and define
\[
T=\{m\in\Z_{>0}:\ \forall p\in\mathbb{P},\ p\mid m\Longrightarrow p\not\equiv 1\pmod 4\}.
\]
Thus the permitted prime divisors are \(2\) and the primes congruent to \(3\pmod 4\). In particular, \(1\in T\), since it has no prime divisors.

Does \(T\) contain infinitely many pairs of consecutive positive integers? Precisely, is it true that
\[
\forall B\in\Z_{>0}\ \exists n\in\Z_{>0}:\quad n\ge B,\qquad n\in T,\qquad n+1\in T?
\]
The supplied record labels this question open. That label is supplied curation, not an independently verified mathematical or literature conclusion. The bibliographic year qualification is recorded below; it does not change the stated mathematical question.
\subsection{Short English statement}
Are there infinitely many positive integers \(n\) such that neither \(n\) nor \(n+1\) has a prime divisor congruent to \(1\pmod 4\)? The prime \(2\) is allowed, and \(1\) satisfies the restriction vacuously.
\subsection{Sources}
\begin{itemize}
\item Local UnsolvedMath registry, supplied record with permanent ID \(30002436\), title ``Consecutive Integers with Restricted Prime Divisors,'' exact source code \texttt{catalogue entry}, in the Oberwolfach Reports Open Problems collection. No individual proposer is supplied; the proposed year is recorded as 2013.
\item Supplied source citation: \emph{Analytic Number Theory (2014)}, Oberwolfach Reports, DOI 10.4171/owr/2013/51, pp.\ 2963--3037. \url{https://doi.org/10.4171/owr/2013/51}. The supplied background describes these reports as records of mathematical workshops and their open-problem discussions. The supplied literature-source entry, labeled ``original OWR report,'' gives this same link.
\item Literature assessment dated exactly \texttt{2026-08-21}: the supplied curation labels the problem open, says that the source poses the consecutive-pair assertion as a question, and reports that no primary proof or disproof was located in its search. This assessment is supplied dated curation, not independently verified; its time-sensitive status and cited claims require independent checking.
\item The supplied citation labels the report 2014, while the proposed year is 2013 and the DOI contains 2013. These bibliographic details are retained as supplied without reconciliation.
\end{itemize}
\end{document}
